\documentclass[12pt,twoside]{article}
\usepackage[portuguese]{babel}
\usepackage[a4paper]{geometry}
\usepackage{graphicx, wrapfig, subcaption, setspace, booktabs}
\usepackage[T1]{fontenc}
\usepackage[utf8]{inputenc}
\usepackage[font=small, labelfont=bf]{caption}
\usepackage{url, lipsum}
\usepackage{hyperref}
\usepackage{xcolor}
\usepackage[portuguese]{babel}
\usepackage{algorithm}
\usepackage{tikz}
\usepackage{indentfirst}
\usepackage{float}
\usepackage[table]{xcolor}
\usepackage[usenames,dvipsnames]{xcolor}
\usepackage[bottom]{footmisc}
\usepackage[document]{ragged2e}
\usepackage{fancyvrb}
\usepackage{caption}
\usepackage{subcaption}
\usepackage{hhline}
\usepackage{circuitikz}
\usepackage{multicol}
\usepackage{colortbl}
\usepackage{xpatch}
\usepackage{varwidth}
\usepackage[cache=false]{minted}
\usepackage{tcolorbox}
\usepackage{etoolbox}
\usepackage{titlesec}
\usepackage{amsfonts}
\usepackage{amssymb}
\usepackage{svg}
\usepackage[neverdecrease]{paralist}
\tcbuselibrary{raster}
\tcbuselibrary{all}
\usetikzlibrary{shadings}
\onehalfspacing
\usepackage{geometry}[]
\geometry{margin=.6in,bottom=1in,top=1in}

\tcbset{skin=enhanced,drop shadow}

\graphicspath{{images/}}

\usepackage{fancyhdr}

\renewcommand{\sectionmark}[1]{\xdef\leftmark{\thesection\quad#1}\xdef\rightmark{}}
\renewcommand{\subsectionmark}[1]{\xdef\rightmark{\thesubsection\quad#1}}

\fancypagestyle{geral}{
  \fancyhf{}
  \renewcommand{\footrulewidth}{0.3pt}
  \fancyhead[R]{\rule[-5mm]{0pt}{2ex}\nouppercase{\leftmark}}
  \fancyfoot[RO, LE]{\nouppercase{\rightmark}}
  \fancyfoot[LO, RE]{\rule[3mm]{0pt}{2ex}\textbf{\thepage}}
}

% \fancypagestyle{firstpage}{
%   \fancyhf{}
%   \renewcommand{\footrulewidth}{0.3pt}
%   \fancyhead[R]{\rule[-5mm]{0pt}{2ex}1\hspace{4mm}Introdução}
%   \fancyfoot[C,RE]{\rule[3mm]{0pt}{2ex}\textbf{\thepage}}
% }

\fancypagestyle{referencias}{
  \fancyhf{}
  \renewcommand{\footrulewidth}{0.3pt}
  \fancyhead[R]{\rule[-5mm]{0pt}{2ex}Referências}
  \fancyfoot[RE, LO]{\rule[3mm]{0pt}{2ex}\textbf{\thepage}}
}


\newcommand{\Sec}[1]{
\section{#1}
\hrule
\vspace{0.5cm}
}

\newcommand{\Subsec}[1]{
\subsection{#1}
\noindent\rule[1.2cm]{3.35in}{.1mm}
\vspace{-0.5cm}
}
%1.675in

\newcommand{\Subsubsec}[1]{
\subsubsection{#1}
}

\makeatletter
\renewcommand\paragraph{\@startsection{paragraph}{4}{\z@}%
            {-2.5ex\@plus -1ex \@minus -.25ex}%
            {1.25ex \@plus .25ex}%
            {\normalfont\normalsize\bfseries}}
\makeatother
\setcounter{secnumdepth}{4} % how many sectioning levels to assign numbers to
\setcounter{tocdepth}{4}    % how many sectioning levels to show in ToC

\newcommand{\Subsubsubsec}[1]{
\paragraph{#1}
}

\newtcolorbox{marker}[1][]{enhanced,
  before skip=2mm,after skip=3mm,
  boxrule=0.4pt,left=5mm,right=2mm,top=1mm,bottom=1mm,
  colback=white,
  colframe=black,
  sharp corners,rounded corners=southeast,arc is angular,arc=3mm,
  underlay={%
    \path[fill=tcbcolback!80!black] ([yshift=3mm]interior.south east)--++(-0.4,-0.1)--++(0.1,-0.2);
    \path[draw=tcbcolframe,shorten <=-0.05mm,shorten >=-0.05mm] ([yshift=3mm]interior.south east)--++(-0.4,-0.1)--++(0.1,-0.2);
    \path[fill=gray!50!black,draw=none] (interior.south west) rectangle node[white]{\Huge\bfseries !} ([xshift=4mm]interior.north west);
    },
  drop fuzzy shadow,#1}
